\chapter{Formation Witnesses}
\label{comp:ch13}

This chapter of the companion mirrors \MT{ch:finite-depth-formation}, with the proofs in
\MT{app:finite-depth-formation-proofs}. It is where the program's central word, formation, receives
theorems, and where the theorems most need to be read against their hypotheses: the headline result
forms collectives out of a hierarchy the model is given, not one it finds. The technical chapter's scope
paragraph declines each stronger reading; the opening paragraph of the introduction summarizes the
chapter more strongly than its theorems support.

\section{The question and the answer}
\label{comp:ch13-question}

Earlier chapters establish that an exact parent exists: coarse-graining the agent network by a Markov
channel gives a normalized parent law that preserves the evidence and splits the free energy exactly
(\MT{thm:rg-exact-coarse-vfe}; see \cref{comp:ch11}). They also establish that this parent is in
general not an autonomous Markov process (\MT{thm:rg-projection-memory}), and how a parent's own
dynamics is fixed by the network it summarizes is the defining open problem of the phase
(\MT{open:rg-pointwise-parent-dynamics}). The formation chapter asks something more concrete. In one
fully normalized experiment at the base point, can coherent groups of agents form one after another
while each parent is started from the complete law of the level below, keeps every loop coordinate and
every attention label, and learns its own copy of a model parameter from its own law? Three statements
answer, and they must be read together.

\MT{thm:finite-depth-adaptive-formation} says yes for a chain of \(2^{L-1}\) triangles at every finite
depth \(L\ge2\): the number of coherent components passes through
\(3\cdot2^{L-1}\to2^{L-1}\to\cdots\to1\) with probability tending to one, and every error made by
replacing a level with its parent vanishes, separately for the free energy, the learning signal, the
attention history, and the continuous path. The succession is put in by hand. The interaction
strengths are supplied as \(\mathrm J_\ell=\epsilon^{-\mathrm Q_\ell}\) with a separation doubly
exponential in the depth, the initial state is a prepared displacement of a common mean, the learned
parameter is pinned by a calibration record, and the attention record is saturated. The chapter says so:
formation in this family is hierarchical relaxation of the prepared mean, and attention and learning
are perturbations the proof controls, not causes (\MT{rem:rg-formation-saturation}). What the theorem
establishes is that exact marked parents can be carried through such a relaxation without their errors
accumulating, at every finite depth, with constants that depend on the depth in an uncontrolled way.

\MT{prop:rg-formation-endogenous-cycle} removes the supplied hierarchy. On a homogeneous cycle of four
agents, where an edge carries energy only while one of its endpoints attends to it, a pair forms by a
feedback between attention and belief relaxation, with probability at least
\(3e^{-1/4}/512\approx4.6\times10^{-3}\) uniformly in the small parameter, and an exact parent follows
the episode for one unit of time. The pairing is seeded by the random initial attention assignment, the
law of the process stays symmetric, there is one level only, and nothing is learned.
\MT{prop:rg-gated-depth-one} shows where this mechanism stops: in the whole gated family on any finite
graph, the coherent blocks of the stationary law are, up to an error vanishing with the temperature, the
components of a minimum edge cover, a forest of stars, and every further merge costs a factor of the
temperature.

The chapter therefore proves that exact parents can track a supplied hierarchy, that one endogenous
level can form on a small graph, and that in the family containing that mechanism the stationary law
stops at one level. It does not prove that collectives form through the agents' coupled inference in
more than one level, the program aim of \MT{sec:intro-program-aim}; that is
\MT{open:rg-formation-general}, and whether learning or loop holonomy can decide which components form
is \MT{open:rg-formation-adaptive-family}. The seeded contrast formation from every Gaussian initial
law, the pairing-selection theorem, and the parent-activation continuation results belong to the
interpretation chapter of the technical manuscript and are treated in \cref{comp:ch14}.

\section{The picture}
\label{comp:ch13-picture}

\paragraph{The triangle chain: springs of very different stiffness.}
Give each agent a belief mean \(\bar k_i\) and a model mean \(\bar m_i\), and each oriented edge
\(e:j\to i\) a translation \(\omega_e\), the link, which shifts the sender's frame into the receiver's.
The mismatch on an edge is \(r_e=\bar k_i-\bar k_j-\omega_e\), and the energy is a sum of squared
mismatches weighted by spring constants. Inside each triangle the springs, and the spring tying each
agent's belief to its own model, have strength \(\mathrm J_1\). The bridge between triangles \(j\) and
\(j+1\) has strength \(\mathrm J_{\ell(j)}\) with \(\ell(j)=2+v_2(j)\), where \(2^{v_2(j)}\) is the largest
power of two dividing \(j\); with four triangles (\(L=3\), twelve agents) the three bridges have levels
two, three, and two. The bridges thus sit on the levels of a binary tree, like the couplings of Dyson's
hierarchical model \citep{Dyson1969}, a resemblance that is structural only. The strengths are
separated by enormous factors, \(\mathrm J_1\gg\mathrm J_2\gg\cdots\gg\mathrm J_L=1\).

The system starts displaced, every mismatch of order \(\mathrm A\), with thermal fluctuations of size at
most \(\sqrt{\mathrm T}=\epsilon^3\mathrm A\), and relaxes by overdamped Langevin dynamics. The stiffest springs
relax first, on times of order \(1/\mathrm J_1\) up to a logarithm, and the triangles become internally
coherent; on times of order \(1/\mathrm J_2\) the level-two bridges relax and triangles pair up; and so on,
until on times of order one the chain is a single component. In the twelve-agent example the counts are
\(12\to4\to2\to1\). An edge is coherent when its conditional root-mean-square mismatch is below
\(\eta_L\mathrm A\), a small fixed fraction of the initial displacement, so formation here means that a
group's internal relative displacement has decayed below a resolution set by the preparation. The mean
dynamics is linear and, relative to the displacement, deterministic up to corrections of order
\(\epsilon^3\), and the grouping is fixed entirely by which springs are strong. A parent is the law with
the already relaxed distortions integrated out; it keeps the loops, the gauge directions, and the
complete attention assignment.

\begin{analogy}{exact}
Linear response after a field is switched off, with adiabatic elimination of the fast modes. The
reference preparation of \MT{def:rg-formation-regime}, a Gaussian with mean \(\mathrm A\mathbf v_0\) and
covariance \(\mathrm TH^{-1}\) on the non-gauge coordinates, is exactly the Gibbs state of
\(\tfrac12\mathbf x^\top H\mathbf x-\mathbf h^\top\mathbf x\) in the static field
\(\mathbf h=\mathrm AH\mathbf v_0\), and the chain's dynamics is the Ornstein--Uhlenbeck relaxation after
the field is switched off at time zero. Integrating out relaxed modes is exact Gaussian integration, with
the Schur complement \(H_\ell\) as retained stiffness. The identification is a theorem for the Gaussian
core, where the record is off and the parameter frozen and the mean follows
\(\bar{\mathbf x}(t)=e^{-H_{\ell'}t}\bar{\mathbf x}(0)\) (\MT{prop:rg-formation-gaussian-core}); in the full
theorem the record and learning add perturbations that the proof bounds.
\end{analogy}

\paragraph{Why the regime is so extreme.}
The reason, informally, is that two requirements pull against each other. Coherence must be many thermal
widths wide, so \(\mathrm J_\ell\mathrm A^2/\mathrm T\) must diverge. But a parent integrates out fast modes
while slower modes are still moving, and the fast modes lag behind the position the slow modes dictate.
The parent's error is that lag measured in thermal units, and at low temperature a small lag is many
thermal widths. In the Gaussian core the loss is exactly of order
\((\mathrm A^2/\mathrm T)(\mathrm J_{\ell+1}^4/\mathrm J_\ell^3)\) (\MTeq{eq:rg-formation-core-window}); keeping
it small while \(\mathrm A^2/\mathrm T\) is large forces each level's stiffness to beat the next by a large
power, and iterating over \(L\) levels gives a separation doubly exponential in \(L\). Separately, the
regime makes the network's stored energy, of order \(\mathrm A^2\mathrm J_1\), negligible against the
unit-scale calibration term. That is what makes the amplitude tiny, and it is what makes learning blind
to the network.

\paragraph{The four-cycle: attention feedback and a frustrated loop.}
Four agents sit on a ring, each receiver attending one of its two neighbors. The prepared state puts
every mean at zero and a unit translation on every link, so every edge has mismatch \(-1\) and the links
carry a circulation of four. If the assignment is a perfect matching, only the two matched edges carry
energy: they contract, while the unattended edges absorb the circulation and stretch. Each receiver's
partner edge soon becomes much cheaper than its exterior edge, and at \(\mathrm T=\epsilon^8\) the attention
posterior then keeps the partner. Attention caused the contraction and the contraction justifies the
attention. Nothing in the energies prefers one pair; the initial assignment, drawn uniformly, decides, and
the law stays invariant under the symmetries of the ring at every time.

\begin{analogy}{structural}
A frustrated plaquette, with symmetry broken along single trajectories. The flux around the loop forbids
all four bonds from being satisfied at once, so the strain must sit somewhere, and attention decides
where. The broken symmetry is of the finite-volume kind: each sample path picks a pair and holds it over
the episode while the ensemble law stays symmetric, as a finite Ising magnet below its ordering
temperature has a symmetric Gibbs state although single trajectories sit in one well. No thermodynamic
limit is taken; law-level selection is the subject of \cref{comp:ch14}.
\end{analogy}

\paragraph{Why the gated mechanism stops at one level.}
On a fixed attention assignment, a sheet, the energy stiffens exactly the attended residuals; every other
direction, including any rigid translation of one block relative to another, is held only by the reference
prior, with stiffness of order \(\mathrm T\) in energy units. All sheets share the minimum energy zero, so a
sheet's stationary weight is its phase-space volume, of order \(\mathrm T^{\mathrm k_S/2}\) with \(\mathrm k_S\)
the number of stiff directions. The law prefers the sheets with the fewest stiff directions, and each
attended edge joining two blocks adds two of them and costs a factor \(\mathrm T\). Blocks float relative to
each other and nothing pulls them together.

\begin{analogy}{structural}
Order by disorder. Among degenerate ground states, harmonic fluctuations favor those with the most soft
modes, because each stiff mode contributes a factor proportional to \(\sqrt{\mathrm T}\) to the partition
function; this is how fluctuations lift a ground-state degeneracy in frustrated magnets. Here the
degenerate states are the attention sheets and the selection favors fragmentation. The counting is the
same; the degenerate set is a finite family of label sectors rather than a continuous manifold.
\end{analogy}

\section{Objects and notation}
\label{comp:ch13-notation}

The technical chapter uses a local notation table (\MT{tab:rg-formation-notation}). The decoder covers
what a first reading needs.

\begin{decoder}
Configuration \(\mathbf x=(\bar k,\bar m,\omega)\) & Belief means, model means, link translations; agent laws are Gaussian with fixed variance, so only means move & Lattice fields: two matter fields on sites, links on bonds \\
Link \(\omega_e\), group \((\R,+)\) & Independent edge variable shifting the sender's frame into the receiver's & Noncompact abelian gauge link \\
Residual \(r_e\), cost \(\tfrac12r_e^2\) & Edge mismatch; the cost is a relative entropy between unit Gaussians & Covariant difference; bond energy \\
Loop \(h_q=\mathfrak h_q/\sqrt3\) & Normalized total translation \(\mathfrak h_q\) around a triangle & Plaquette flux \\
\(\mathcal U,\mathcal H,\mathcal K\) & Normal (gauge-invariant, flux-free) distortions; loops; pure-gauge directions & Massive modes; flux sector; gauge orbit \\
Stiffness \(H\); strengths \(\mathrm J_\ell=\epsilon^{-\mathrm Q_\ell}\) & Quadratic form of the energy; supplied spring constants, one per level & Dynamical matrix; coupling hierarchy \\
\(\mathrm A\), \(\mathbf v_0\), margin \(\mathrm m_L\) & Size and direction of the initial displacement; unrelaxed bridges stay at least \(\mathrm m_L\mathrm A/2\) from zero & Applied field switched off at time zero \\
Calibration record \(o_{\mathrm{cal}}=2\) & Precise observation of the learned parameter itself & External clamp \\
Learned scalar \(\vartheta\) & Global multiplier of the stiffness, \(\vartheta H\), learned by natural gradient & Slowly ramped coupling \\
Source, assignment \(S\), sheet & Which neighbor each receiver attends; one fixed assignment and its slice of state space & Label field; sector \\
Source record \(E_i\), \(\phi_S\), \(B_0\) & Bounded observation of the attended edge's mismatch, total effect at most \(B_0=10^{-4}\) & Weak bounded potential \\
\(\boldsymbol\Pi\), \(q^\star\), \(q_t\) & Physical posterior; best law of the restricted family; evolving recognition law & Gibbs state; variational optimum; time-dependent density \\
Marked channel \(C_{\ell\leftarrow\ell'}\), parent & Keep retained coordinates and every label, integrate out the rest & Exact decimation \\
Schur complement \(H_\ell\) & Stiffness of retained coordinates after minimizing over discarded ones & Tree-level effective stiffness \\
Discarded loss \(\Delta_{\ell':\ell}\) & Mean relative entropy of the discarded coordinates' actual law to their equilibrium law, given the retained ones & Adiabatic-elimination error \\
Window \(I_\ell\), activation, \(k_L\) & Interval on which level \(\ell\) is coherent; start of the next parent; how many relaxation times to wait & Timing of decimations \\
Coherent edge, component & Conditional root-mean-square residual at most \(\eta_L\mathrm A\); cluster of such edges & Cluster at a resolution \\
\end{decoder}

Some traps. \(\mathrm T\) is the temperature, not the frame comparison \(T_{ij}\). \(\mathrm J_\ell\) is a spring
constant while calligraphic \(\mathcal J\) is the set of assignments, and a source is the neighbor a
receiver attends, not a source term in a field theory. \(H\) is a stiffness matrix, not a Hamiltonian or a
holonomy. The root is the set of gauge directions, not the root of a tree. A parent is a marginal law,
not an agent: the scope paragraph declines that parents are coarse agents with their own belief and model
fibers. Learning is the deterministic drift of one scalar. Coherence is measured against the preparation
amplitude \(\mathrm A\), not the thermal width. Complete, exact, marked, and retained are emphasis on
objects that keep all their information; declared marks a modeling choice, as in the declared parent
kinetics. NOT-CLAIMED marks a declined claim, and declining is not refuting.

\section{The formulas, explained}
\label{comp:ch13-formulas}

\subsection{The carrier}
\label{comp:ch13-carrier}

The chain of \MT{def:rg-formation-carrier} has \(N=3\cdot2^{L-1}\) agents and configuration space
\(\R^{10\cdot2^{L-1}-1}\). By \MT{lem:rg-formation-carrier},
\[
\mathbf x^\top H\mathbf x=\mathrm J_1\sum_i\mathfrak d_i(\mathbf x)^2+\sum_e\mathrm J_{\ell(e)}r_e(\mathbf x)^2,
\qquad
|r_q|^2=|O_\triangle r_q|^2+h_q^2 .
\]
The first sum ties each belief to its own model through the discrepancy
\(\mathfrak d_i=\bar k_i-\bar m_i\), with the strongest spring; the second weights each edge mismatch by its
level. A passive frame change \(\bar k_i\mapsto\bar k_i+g_i\), \(\bar m_i\mapsto\bar m_i+g_i\),
\(\omega_e\mapsto\omega_e+g_i-g_j\) leaves every \(r_e\) and \(\mathfrak d_i\) unchanged, the
translation-group case of the frame invariance proved in general
(\MT{cor:elbo-independent-frame-invariance}). The second identity splits a triangle's squared mismatches
into two gauge-invariant contrasts and the flux \(h_q^2\): the vertex means telescope around the triangle,
so the mismatches sum to \(-\mathfrak h_q\) whatever the means are, and a triangle with flux cannot be
fully satisfied. Space splits orthogonally into \(\mathcal U\oplus\mathcal H\oplus\mathcal K\), and \(H\)
vanishes exactly on \(\mathcal K\).

\begin{analogy}{exact}
The quadratic form is the action of a Gaussian, noncompact lattice abelian gauge-Higgs model in
St\"uckelberg form, in the lattice formulation of \citet{WilsonConfinement1974} but with no plaquette
term: flux costs energy only through the matter coupling, and the model channel is a second matter field
tied on site to the first. The group is \((\R,+)\), abelian and one-dimensional, so nothing about
nonabelian or compact groups can be read off.
\end{analogy}

\subsection{One normalized experiment and its three laws}
\label{comp:ch13-experiment}

\MT{def:rg-formation-experiment} makes the carrier a probability model. Each agent has a model sample
\(m_i\sim\mathcal N(\bar m_i,\sigma^2_{\mathrm{ev}})\) and, given it, a belief sample
\(k_i\sim\mathcal N(m_i,1-\sigma^2_{\mathrm{ev}})\). The recognition family fixes the coupling of the two
samples to its optimal form, at the price \(\mathfrak d_i^2/(2(1-\sigma^2_{\mathrm{ev}}))\) per agent
(\MT{prop:rg-formation-normalized-experiment}), summing to \(\tfrac12\mathbf x^\top R_{\mathrm{ev}}\mathbf x\).
The configuration prior is Gaussian with precision \(\Omega_0(\vartheta)/\mathrm T\), where
\[
\Omega_0(\vartheta)=\vartheta H+P_{\mathcal K}-\mathrm TR_{\mathrm{ev}} .
\]
The first term is the physics, a stiffness multiplied by the learned
\(\vartheta\); the second normalizes the gauge directions; the third cancels the restriction price so that
the optimizer has a temperature-independent energy. The last two are bookkeeping.

Three observations complete the joint. A root observation gives the gauge coordinates an independent
Gaussian law, a form of gauge fixing. A calibration record \(o_{\mathrm{cal}}=2\) from
\(\mathcal N(\vartheta,\mathrm T)\) observes the learned parameter itself; it is what the parameter will
learn. And each receiver has a binary record reading the edge it attends,
\[
\Pr(E_i=1\mid\mathbf x,S_i=j)=\tfrac12\exp\Bigl[-\delta\chi\Bigl(\sqrt{\mathrm J_1/\mathrm T}r_{ij}(\mathbf x)\Bigr)\Bigr],
\qquad \chi(r)=1-e^{-r^2/2},\qquad \delta=10^{-4}/N .
\]
The argument of \(\chi\) is the mismatch in units of the thermal width \(\sqrt{\mathrm T/\mathrm J_1}\), and
\(\chi\) saturates at one. Summed over receivers, the configuration-dependent part \(\phi_S\) of the record's
negative log-likelihood lies in \([0,10^{-4}]\), so attention differs from its prior by a factor of at most \(e^{\pm10^{-4}}\)
(\MTeq{eq:rg-formation-source-saturation}): attention is blind to the configuration by construction. The
attention prior is uniform except at one receiver \(i_0\), whose prior is
\(\operatorname{sig}(4(\vartheta-3/2))\); the learned parameter controls one attention weight, through the
prior only.

Three laws coexist (\MT{def:rg-formation-laws}; \MT{prop:rg-formation-laws}): the physical posterior
\(\boldsymbol\Pi\), the restricted optimizer \(q^\star\), and the evolving recognition law \(q_t\). The
scaled free energy is
\(\Fenergy_{\mathrm T}(q,\vartheta)=\mathrm T\KL(q\Vert q^\star_\vartheta)+\mathrm C_{\mathrm T}(\vartheta)\),
and the optimizer differs from the posterior by at most
\(2e^{B_0}N\mathrm T^2/((1-\sigma^2_{\mathrm{ev}})\mathrm J_1)\) in free energy
(\MTeq{eq:rg-formation-physical-optimizer-gap}), negligible here. The theorem is proved for both
\(\boldsymbol\Pi\) and \(q^\star\).

\subsection{Marked parents and the discarded loss}
\label{comp:ch13-parents}

The parent at level \(\ell\) comes from the deterministic channel
\(C_{\ell\leftarrow\ell'}((y,S),\cdot)=\delta_{(P_\ell y,S)}\), which projects onto the retained space
\(W_\ell\) (unrelaxed normal directions, all loops, all gauge directions), keeps the whole assignment, and
leaves the observations alone (\MT{def:rg-formation-parents}). The parent law is a marginal, with sheet
potential
\[
V_{\ell,S}(y,\vartheta)=\tfrac12\vartheta y^\top H_\ell y+\mathrm C_\ell^{\mathrm G}(\vartheta)+\mathrm T\psi_{\ell,S}(y,\vartheta),
\qquad 0\le\psi_{\ell,S}\le B_0,
\]
by \MT{prop:rg-formation-parent-potential}: the Schur-complement stiffness, the Gaussian determinant of
the integrated modes (their entropy, which depends on \(\vartheta\)), and the record averaged over the
discarded coordinates, still below \(10^{-4}\). The chain rule of relative entropy gives what a parent
loses,
\[
\Fenergy^{\ell'}_{\mathrm T}(q^{\ell'},\vartheta)-\Fenergy^{\ell}_{\mathrm T}\bigl((C_{\ell\leftarrow\ell'})_\#q^{\ell'},\vartheta\bigr)
=\mathrm T\Delta_{\ell':\ell}(q^{\ell'},\vartheta),
\]
where \(\Delta_{\ell':\ell}\) is the mean relative entropy of the discarded coordinates' actual conditional
law to the optimizer's conditional. It vanishes exactly when the discarded coordinates are in conditional
equilibrium, the condition for exact adiabatic elimination. These identities are standard in substance.

The parent is not a coarse agent. At the top level the retained space has dimension \(4\cdot2^{L-1}\): all
\(N\) gauge directions and all \(N/3\) loops, with every normal direction integrated out
(\MTeq{eq:rg-formation-dimensions}), alongside the full assignment with \(2^N\) values. For \(L=2\) the
carriers go \(\R^{19}\to\R^9\to\R^8\) with 64 assignments throughout, and the hidden conditionals are not
Gaussian once the record is on (\MT{rem:rg-formation-two-triangle}). The top parent is the zero-mode and
flux sector of the Gaussian theory together with the complete label field.

\subsection{Kinetics and learning}
\label{comp:ch13-kinetics}

The microscopic kinetics (\MT{def:rg-formation-kinetics}) are overdamped Langevin dynamics on each sheet
with a rate-one clock that redraws the whole assignment from the attention posterior,
\[
\mathscr L_\vartheta^0\mathsf f(\mathbf x,S)=\mathrm T\Delta_{\mathbf x}\mathsf f-\nabla_{\mathbf x}U_{\vartheta,S}\cdot\nabla_{\mathbf x}\mathsf f
+\sum_{S'}\beta_{S'}(\mathbf x,\vartheta)\bigl[\mathsf f(\mathbf x,S')-\mathsf f(\mathbf x,S)\bigr],
\]
a diffusion at temperature \(\mathrm T\) with unit mobility plus a heat-bath update of the labels. The parent
generator has the same form with the parent potential and inherited jump rows, conditional averages of
the fine rows, which are reversible and compose across levels (\MT{prop:rg-formation-rate-composition}).
The parent kinetics is declared, not derived: the time evolution of a true marginal carries memory
(\MT{thm:rg-projection-memory}; compare \citet{Zwanzig1960}), and the chapter checks its Markov closure
afterwards through conclusions (c) and (e).

Each level learns its own copy of \(\vartheta\) by natural gradient,
\[
\dot\vartheta_\ell=-\epsilon\mathsf m(\vartheta_\ell)^{-1}g_\ell(q^\ell_t,\vartheta_\ell),
\qquad
\mathsf M(\vartheta)=\frac1{\mathrm T}+\frac12\Tr\bigl[(\Omega_0(\vartheta)^{-1}H)^2\bigr]+16\pi_{i_0,+}(1-\pi_{i_0,+}),
\]
with \(\mathsf m=\mathrm T\mathsf M\) (\MT{def:rg-formation-learning};
\MT{prop:rg-formation-score-dissipation}). The metric is dominated by the calibration term, so
\(\mathsf m\approx1\), and the covector \(g_\ell\) (\MTeq{eq:rg-formation-complete-score}) is \(\vartheta-2\)
plus the stored elastic energy \(\tfrac12\E y^\top H_\ell y\), at most of order \(\mathrm A^2\mathrm J_1\),
plus entropic and attention terms of order \(\mathrm T\). Hence
\(\vartheta(t)\approx2-e^{-\epsilon t}\). Along the coupled dynamics the free energy decreases through
three nonnegative dissipation terms, the Fisher information of the diffusion, the symmetrized entropy
production of the label jumps, and the natural-gradient term \(\epsilon\mathsf m^{-1}g_\ell^2\)
(\MTeq{eq:rg-formation-dissipation}): an H-theorem, standard in substance.

\subsection{The regime}
\label{comp:ch13-regime}

The regime of \MT{def:rg-formation-regime} is
\[
\mathrm Q_L=0,\quad \mathrm Q_\ell=2\mathrm Q_{\ell+1}+20,\quad \mathrm J_\ell=\epsilon^{-\mathrm Q_\ell},\quad
\mathrm A^2=\epsilon^{3\mathrm Q_1+20},\quad \mathrm T=\mathrm A^2\epsilon^6,\quad t_*=\frac2\epsilon .
\]
Then \(\mathrm Q_1=20(2^{L-1}-1)\) and \(\mathrm T=\epsilon^{60\cdot2^{L-1}-34}\) (\MT{rem:rg-formation-regime}).
Since \(N=3\cdot2^{L-1}\), this is \(\mathrm T=\epsilon^{20N-34}\): the temperature exponent is linear in the
number of agents, and the stiffness range is \(\epsilon^{-20(N/3-1)}\). The last two columns of the table
are arithmetic on the class of schedules in which the sharpness problem is posed
(\cref{comp:ch13-core}): the least admissible \(\mathrm Q_1\) there, rounded, and its temperature.

\begin{center}
\small
\begin{tabular}{@{}rrrrrrr@{}}
\toprule
\(L\) & \(N\) & \(\mathrm Q_1\) & \(\mathrm A\) & \(\mathrm T\) & least \(\mathrm Q_1\), base \(4/3\) & its \(\mathrm T\)\\
\midrule
2 & 6 & 20 & \(\epsilon^{40}\) & \(\epsilon^{86}\) & 14 & \(\epsilon^{68}\)\\
3 & 12 & 60 & \(\epsilon^{100}\) & \(\epsilon^{206}\) & 25.3 & \(\epsilon^{102}\)\\
4 & 24 & 140 & \(\epsilon^{220}\) & \(\epsilon^{446}\) & 40.4 & \(\epsilon^{147}\)\\
8 & 384 & 2540 & \(\epsilon^{3820}\) & \(\epsilon^{7646}\) & 171 & \(\epsilon^{539}\)\\
9 & 768 & 5100 & \(\epsilon^{7660}\) & \(\epsilon^{15326}\) & 235 & \(\epsilon^{730}\)\\
\bottomrule
\end{tabular}
\end{center}

The ratios the proof actually uses are the stiffness range \(\epsilon^{-\mathrm Q_1}\), the noise-to-signal
ratio \(\mathrm T/\mathrm A^2=\epsilon^6\), the learning mobility \(\epsilon\), and the stored energy relative to
calibration, \(\mathrm A^2\mathrm J_1=\epsilon^{2\mathrm Q_1+20}\). The window of level \(\ell\) is
\(I_\ell=[t_\ell,2t_\ell]\) with \(t_\ell=k_L\log(1/\epsilon)/\mathrm J_\ell\); the last closes at order
\(k_L\log(1/\epsilon)\), long before the horizon \(2/\epsilon\) on which the parameter moves. The entrance
constant must exceed thresholds proportional to \(\mathrm Q_1\) (\MTeq{eq:rg-formation-common-entrance}), and
\(\epsilon_L\) is not computed. The regime is a formal asymptotic regime at each fixed depth, with no limit in
which \(N\) grows at fixed \(\epsilon\).

\subsection{The theorem, conclusion by conclusion}
\label{comp:ch13-theorem}

\MT{thm:finite-depth-adaptive-formation} assumes the experiment and kinetics at a fixed depth (H1), the
regime (H2), an initial law in the preparation neighborhood (H3), and scheduled activations (H4). The
preparation neighborhood consists of smooth transports of the Gaussian reference that move each sample by
less than \(\epsilon\sqrt{\mathrm T}\) in the energy norm (\MTeq{eq:rg-formation-preparation-transport}): a
thermalized cloud around the displaced mean, slightly deformed, possibly depending on \(\epsilon\) and
\(L\). The reader's phrase ``coherence is linear hierarchical relaxation of the prepared mean'' has a
precise meaning here. Up to perturbations the proof bounds, the mean obeys
\(\dot{\bar{\mathbf x}}=-\vartheta H\bar{\mathbf x}\) from \(\mathrm A\mathbf v_0\). In the flag coordinates of
\MTeq{eq:rg-formation-weighted-flag}, layer \(\ell\) has stiffness of order \(\mathrm J_\ell\) and couples to the
slower layers only through strengths of order \(\mathrm J_{\ell+1}\), so the solution decays level by level. On \(I_\ell\) the layers up to \(\ell\) have decayed
below \(\eta_L\mathrm A\), while the higher bridges have moved by at most of order
\(\mathrm J_{\ell+1}t_\ell\mathrm A\to0\) and keep their prepared margin.

Conclusion (a) gives unique non-exploding solutions for every process, and learned parameters in
\([1,21/10]\) that stay within \(\mathrm C_L\epsilon^{2\mathrm Q_1+20}\) of the calibration curve
\(2-e^{-\epsilon t}\) (\MTeq{eq:rg-formation-parameter-bound}): learning is relaxation to the calibrated value,
the same at every level because no level can see the network.

Conclusion (b) is the formation statement. Except on an event of probability at most
\(\mathrm C_L(1+t_*)\exp(-\mathrm c_L/\epsilon^6)\), the fine process and its parents show the counts
\(3\cdot2^{L-1}\to2^{L-1}\to\cdots\to1\) (\MTeq{eq:rg-formation-agent-counts}): on \(I_\ell\) every edge of
level at most \(\ell\) is coherent, every belief agrees with its model to within \(\eta_L\mathrm A\), and every
higher bridge stays a fixed multiple of \(\mathrm A\) from zero; the final component persists to the horizon,
and the loops stay random and retained. The exponent \(\epsilon^{-6}=\mathrm A^2/\mathrm T\) is the
signal-to-noise ratio. By \MT{rem:rg-formation-saturation} the coherence proof uses the record only through
\(0\le\phi_S\le B_0\) and the parameter only through \(1\le\vartheta\le21/10\), so it gives the same
components with the record off and the parameter frozen at one. Attention and learning play no role in (b).

Conclusion (c) bounds the free energy discarded at each activation, \(\Delta_{\ell-1:\ell}\le\mathrm C_L\epsilon^{14}\)
early and \(\mathrm C_L\epsilon^2\) at the last level, and the difference of consecutive learning signals by
\(\mathrm C_L\mathrm T\epsilon^7\) and \(\mathrm C_L\mathrm T\epsilon\) (\MTeq{eq:rg-formation-conditional-score-bounds}):
the parents are faithful in free energy and in learning gradient. The exponent 14 is
\(\mathrm A^2\mathrm J_{\ell+1}^2/(\mathrm T\mathrm J_\ell)=\epsilon^{-6}\epsilon^{20}\), the general Fisher-budget
estimate; the Gaussian core shows the true early loss is much smaller.

Conclusion (d) shows the model is not static: from every activation to the horizon the distinguished
receiver's mass on source \(+\) rises by more than one half, and the first hidden conditionals change by
relative entropy more than \(1/200\) (\MTeq{eq:rg-formation-source-movement}). Both are driven by
\(\vartheta\) moving toward two, through the prior \(\operatorname{sig}(4(\vartheta-3/2))\) and the stiffness
\(\vartheta H\), and both happen after formation, since \(\vartheta\le1+\mathrm C_L\epsilon\log(1/\epsilon)\)
until the last window closes. This is an engineered stress test showing that the lineage tracks kernels
that really move, not a mechanism of formation.

Conclusion (e) says that after each activation the retained coordinates of the fine process and of the
parent stay within \(\mathrm C_L\mathrm A\epsilon^2\) of one deterministic path except with probability
\(\mathrm C_L(1+t_*)e^{-\mathrm c_L/\epsilon^2}\) (\MTeq{eq:rg-formation-common-path-concentration}), and that
their attention histories can be coupled to disagree with probability at most \(\mathrm C_L\epsilon^6\)
(\MTeq{eq:rg-formation-two-path-bounds}). The resolution \(\mathrm A\epsilon^2=\sqrt{\mathrm T}/\epsilon\) lies
between the thermal width and the displacement. This small-noise statement is weaker than closeness of
path laws, which the chapter declines along with exact lumpability: the parent reproduces the fine
trajectory at a resolution coarser than the fluctuations it carries.

\MT{rem:rg-formation-saturation} fixes the reading of the whole theorem. The record is saturated, since at a
residual of order \(\mathrm A\) its argument is of order \(\epsilon^{-(3+\mathrm Q_1/2)}\), so it cannot tell a
coherent edge from an incoherent one at the coherence resolution; learning is calibration; and the loops
have root-mean-square size of order \(\sqrt{\mathrm T/\mathrm J_1}\), so holonomy is thermal noise. The
chapter's own summary says the reductions succeed in a regime ``in which nothing adaptive decides the
outcome''.

\subsection{The Gaussian core and whether the doubling is needed}
\label{comp:ch13-core}

\MT{prop:rg-formation-gaussian-core} computes the loss with the record off and the parameter frozen. Part (a)
gives \(\Delta=\tfrac1{2\mathrm T}\mathfrak v^\top H_{zz}^{-1}\mathfrak v\), with \(\mathfrak v\) the force on the
discarded coordinates: the lag of the fast modes behind their slaved position, weighted by compliance, in
thermal units. Part (b) gives, uniformly on each early window,
\[
\Delta_{\ell-1:\ell}(t)=\frac{\mathrm A^2}{\mathrm T}\frac{\mathrm J_{\ell+1}^4}{\mathrm J_\ell^3}
\bigl[\mathfrak y_\ell(\mathbf v_0)+o(1)\bigr],
\qquad \mathfrak y_\ell(\mathbf v_0)>0 .
\]
Informally, the next level moves at rate \(\mathrm J_{\ell+1}\) and pulls through couplings of order
\(\mathrm J_{\ell+1}\), the current level answers with compliance \(1/\mathrm J_\ell\), so the lag force is of
order \(\mathrm A\mathrm J_{\ell+1}^2/\mathrm J_\ell\) and its compliance-weighted square gives the display;
the general estimate behind (c) is larger by \((\mathrm J_\ell/\mathrm J_{\ell+1})^2\). Part (c) turns this into
a criterion: with \(\mathrm A^2/\mathrm T=\epsilon^{-\mathrm Q_{\mathrm T}}\) and a large enough entrance constant,
the core loss vanishes on every early window exactly when \(3\mathrm Q_\ell-4\mathrm Q_{\ell+1}>\mathrm Q_{\mathrm T}\) for every \(\ell<L\), and is
\(O(\epsilon^{\mathrm Q_\Delta})\) exactly when
\(3\mathrm Q_\ell-4\mathrm Q_{\ell+1}\ge\mathrm Q_{\mathrm T}+\mathrm Q_\Delta\), forcing
\(\mathrm Q_1\ge(\mathrm Q_{\mathrm T}+\mathrm Q_\Delta)[(4/3)^{L-1}-1]\).

Is the doubling required? Not shown, and the chapter says so. It enters the proof through the full Fisher
budget of \MT{lem:formation-conditional-entropy} and the history estimate of
\MT{lem:formation-history-coupling}. What is proved is that in the Gaussian core, for this preparation and
schedule, a separation growing like \((4/3)^L\) in the exponent is necessary. Whether the full family works
under the base-\(4/3\) schedules \(3\mathrm Q_\ell\ge4\mathrm Q_{\ell+1}+20\), \(\mathrm Q_\ell\ge14\), with
depth-independent constants, is \MT{open:rg-formation-regime-sharpness}. Either way the separation diverges
with the number of agents: since \(L-1=\log_2(N/3)\), \((4/3)^{L-1}=(N/3)^{\log_2(4/3)}\approx(N/3)^{0.415}\), so
the exponent of the stiffness range grows linearly in \(N\) under the doubling and like a power of \(N\) under
the best schedule the core allows. That is the quantitative sense in which this family has no large-\(N\)
limit.

\subsection{The four-cycle witness}
\label{comp:ch13-cycle}

\MT{def:rg-formation-cycle-experiment} puts four agents on a ring with edge cost and record
\[
\mathcal D_{e_i}(\mathbf x)=\tfrac12\bigl[(r^b_i)^2+(r^m_i)^2\bigr],
\qquad
\Pr(E_i=1\mid\mathbf x,S)=\tfrac12\exp\bigl[-\mathcal D_{e(i,S_i)}(\mathbf x)/\mathrm T\bigr],
\]
the cost adding belief- and model-channel relative entropies and the record reading only the attended
edge, with weight \(1/\mathrm T\). An attended edge is stiff against the thermal scale; an unattended edge
carries no energy at all. This is the gating. The prior is uniform, nothing is learned, every coordinate
has unit mobility, \(\mathrm T=\epsilon^8\) with \(\epsilon\le1/100\), the formation time is
\(t_\epsilon=2\log(1/\epsilon)\), and the horizon is \(t_\epsilon+1\). The preparation is an equal mixture of
narrow Gaussians around \(\pm\mathbf x^\circ\), where \(\mathbf x^\circ\) has all means zero and all links one,
with a uniform assignment, so it names no pair. On the path that holds a matching from \(\mathbf x^\circ\),
\(r_{\mathrm{int}}=-e^{-(8+\mathrm T)t}\) and \(r_{\mathrm{ext}}=e^{-\mathrm Tt}[-\tfrac32+\tfrac12e^{-8t}]\)
(\MTeq{eq:rg-formation-cycle-center}): matched mismatches decay at rate eight, exterior ones go to about
\(-3/2\) and then decay at the thermal rate. The circulation, equal to minus the sum of the belief mismatches,
drops from four to about three and then relaxes only at rate \(\mathrm T\); it is what stretches the
unattended edges.

Part (a) of the proposition gives the episode: with probability at least \(p_0=3e^{-1/4}/512\) the assignment
is one perfect matching throughout; at time zero each matched pair has belief relative entropy above
\(1/4\), so the contrast is entered during the episode; after \(t_\epsilon\) the matched relative entropies are
at most \(\epsilon^4\); and after time \(1/4\) every exterior belief relative entropy exceeds \(1/2\). The pair
is held with restoring rates \(8+\mathrm T\) and \(4+5\mathrm T\), and the kinetics started from the
stationary law realize the event with probability at most \(4e^{-1/(4\mathrm T)}<p_0/2\). Part (b) is the
parent channel, discarding the three fast normal directions of the attended edges and keeping the
assignment, the loop, and one fast coordinate. Part (c) bounds the discarded conditional relative entropy
by \(\bar\eta_\epsilon\) (\MTeq{eq:rg-formation-cycle-constants}). Part (d) puts the parent activated at
\(t_\epsilon\) within total variation \(2p_\epsilon/(1-p_\epsilon)+4\sqrt{8/3}e^{-1/\mathrm T}\) of the native
retained path law conditioned on the same prefix, assignment path included. The constants are explicit at every \(\epsilon\le1/100\); at
\(\epsilon=1/100\) they give a path-law error of about \(1.5\) percent and a discarded loss of about \(0.045\)
nats, both vanishing as \(\epsilon\to0\). Unlike the chain, the comparison here is in total variation of path
laws, because on a held matching the hidden coordinates are exactly independent Ornstein--Uhlenbeck modes.

\begin{center}
\small
\begin{tabular}{@{}>{\raggedright\arraybackslash}p{0.26\textwidth}>{\raggedright\arraybackslash}p{0.54\textwidth}>{\raggedright\arraybackslash}p{0.14\textwidth}@{}}
\toprule
Ingredient & Role & Origin \\
\midrule
Attention-contraction feedback & Attended edges contract, contracted edges win attention, the matching persists & endogenous \\
Energy scale & All edges have the same cost structure & homogeneous \\
Which pair forms & Initial assignment, drawn uniformly; the law stays symmetric & random seed \\
Circulated preparation & A flux around the ring that survives on the episode's time scale and stretches the unattended edges & supplied \\
Source gating & Energy only while attended, with weight \(1/\mathrm T\) & supplied \\
Temperature \(\epsilon^8\) & Makes attention an arg-min of edge costs & supplied \\
Euclidean kinetic metric & Beliefs, models, and links all move at rate one; models are not slow & supplied \\
Duration-one episode; no learning & The pair is proved held for one time unit; parameters fixed & supplied \\
\bottomrule
\end{tabular}
\end{center}

One observation places the witness. The four-cycle belongs to the gated family, so by
\MT{prop:rg-gated-depth-one} its stationary law already sits on the two perfect-matching sheets, up to a
fraction of order \(\mathrm T\) of its mass, with the matched pairs as its coherent blocks. The paired block
partition the episode reaches is the typical partition of equilibrium; what distinguishes the episode from
equilibrium coherence is its entrance clause (the pairs are incoherent at time zero) and the seeded choice
of matching. Physically, the witness is relaxation from
a prepared symmetric incoherent start into one of two paired equilibrium sectors, the choice made by a
sampled label and held by attention feedback. That is a genuine coupled-inference mechanism at one level,
not a collective the equilibrium law would not already contain.

\subsection{The gated family and depth one}
\label{comp:ch13-gated}

\MT{def:rg-gated-translation-family} places the four-cycle's laws, records, and kinetics on any finite
simple graph without isolated vertices, with a product prior independent of \(\mathrm T\). On a sheet the
energy is \(V_S=\tfrac12\mathbf x^\top H_S\mathbf x\), where \(H_S\) sums over receivers the residual pair of
the attended edge, and both laws are Gaussian on each sheet with precision \(H_S/\mathrm T+\Omega\)
(\MTeq{eq:rg-gated-sheet-law}). By \MT{prop:rg-gated-depth-one}, (a) every sheet has a stiff band of order
one on the \(\mathrm k_S\) attended directions and a soft band of order \(\mathrm T\) elsewhere, and the prior
enters only sheet weights, never stiffness; (b) the configuration marginal is a Gibbs law within
\(\mathrm T\sum_i\log(1/\min_j\pi_{ij})\) of the energy \(V_{\min}=\sum_i\min_j\mathcal D_{e(i,j)}\), in which each
receiver pays only for its cheapest edge and which is flat under rigid translation of any block of a
covering edge set; (c) the counting is
\[
\operatorname{rank}\mathbf a_F=|F|+N-\operatorname{cc}(F),
\qquad
\mathrm k_S\ge2(N-\mathfrak n),
\qquad
q_\Omega\{S:\operatorname{cc}(\operatorname{att}(S))\le\mathfrak n-m\}\le\mathrm K_m\mathrm T^m,
\]
with equality in the middle exactly for minimum edge covers, the spanning forests of \(\mathfrak n\) stars,
and sheet weights proportional to \(\mathrm T^{\mathrm k_S/2}\) up to controlled factors
(\MTeq{eq:rg-gated-sheet-weight}); (d) at resolution \(\mathrm a\) the stationary block partition differs from
a minimum-edge-cover partition with probability at most
\[
\mathrm K_1\mathrm T+\frac{|E|\|\Omega\|}{2\mathrm s_*^2}\mathrm a^2+4|E|e^{-\mathrm a^2/(4\mathrm T)},
\]
the mass of merged sheets, the chance that an unheld edge between blocks falls inside the resolution, and
the chance that a held edge fluctuates outside it. On an even ring or an even complete graph the blocks are
the pairs of a perfect matching, and on a star graph the block is the star: the stationary law fragments
into the smallest blocks the graph allows.

Is this large-\(N\), and is it static? It holds on every finite graph, so it applies verbatim at 512 or more
agents and is not a small-network result. But it concerns \(\mathrm T\to0\) at a fixed graph with
graph-dependent constants. \(\mathrm K_1\) contains the factor \((1+\|\Omega\|/\mathrm h_*)^{N-\mathfrak n}\),
exponential in \(N\), times a ratio whose numerator runs over every merged assignment
(\MTeq{eq:rg-gated-merge-cost}), and the second term grows with the number of edges. Nothing keeps the bound
small as \(N\) grows at fixed temperature, so the limit of many agents at fixed temperature, where the number
of ways to merge competes with the cost of each merge, is not addressed. The statement is also purely
static: it has nothing on relaxation times, metastability, or how long sampled paths can hold a larger block.
Since the dynamic question is posed only outside the family (\MT{open:rg-formation-general}), a transient
second coalescence inside the gated family is neither excluded nor posed.

\subsection{The authoring rules R13 to R16 in physics terms}
\label{comp:ch13-rules}

Four authoring rules govern what the technical documents may call learning-driven selection, recursion, a
coherence episode, and a settled open problem.

R13, the Curie rule, is Curie's principle applied to learning. If model, records, parameterization, and
kinetics share a symmetry, a claim that learning selects a structure the symmetry does not fix must name
the symmetry-breaking datum it amplifies and compare against a frozen-learning run, or else declare the
selecting mechanism, for example a rule acting on sampled paths, and then the claim concerns sampled paths.
A unique, equivariant law-level dynamics from a symmetric start cannot produce an asymmetric law, the
finite-volume fact that a finite system's Gibbs state is symmetric. The rule is substantive: it is why the
four-cycle result claims only a sampled-path pairing.

R14 defines recursion by a counterfactual. One declares a handoff time, a first-stage event with its
blocks, a second-stage event whose block is a union of first-stage blocks, a statistic such as the second
stage's probability, and a class of handoff states that includes, with probability bounded below, states
from histories without first-stage formation. The second stage is recursive when it has probability
bounded below and its statistic after actual first-stage formation beats its value after a counterfactual
handoff by a margin. If it happens with probability tending to one from every handoff in the class, or its
proof uses only properties that counterfactual handoffs share, the sequence is a supplied-scale sequence:
the hierarchy is in the Hamiltonian, as in any system whose separated couplings order each on its own scale
whether or not the previous scale has ordered. The physical distinction is between a hierarchy of
couplings and a hierarchy of products, whose stage-two units are made of stage-one units and would not
form otherwise. The definition is substantive. The chapter does not run the test on the triangle chain, but
each level's coherence follows from that level's stiffness acting on the prepared displacement, so in
R14's terms the chain is a supplied-scale sequence; this is the companion's reading.

R15 defines a coherence episode. The restoring clause requires a named interaction whose rate times the
squared resolution, divided by the temperature, diverges on the sheets the kinetics select: by
equipartition a mode of stiffness \(k\) has width \(\sqrt{\mathrm T/k}\), so the block sits many thermal widths
inside its resolution, held by energy rather than luck. The horizon clause requires horizons short compared
with the time equilibrium takes to revisit the event, or a probability exceeding the equilibrium value by a
margin. Restored coherence that fails the horizon clause, for instance because equilibrium has it with
probability tending to one, is restored equilibrium coherence, not formation; coherence failing the
restoring clause is coincidental. The recurrence branch is how a physicist says a broken-symmetry state is
real, since tunneling times diverge; the margin branch admits the transients of relaxation from any
prepared, atypical start. The rule thus excludes thermal accidents and equilibrium ordering from the word
formation but not preparation-driven relaxation. The four-cycle meets both clauses, the horizon clause
through its entrance requirement. In the chain the intermediate windows meet the margin branch, but the
final single component is what the restricted optimizer looks like, every residual being thermal under it,
so by R15's own terms the chain's last stage is restored equilibrium coherence. The chapter does not
evaluate R15 for the chain; this is the companion's reading.

R16 requires each open problem to be posed in a stated class with two nontrivial settling branches, neither
met by an existing statement or a trivial construction. It is bookkeeping for well-posedness, and it is why
the open problems below are long and cannot be closed by pointing at the witnesses already proved.

\section{What is proved, and how solid it is}
\label{comp:ch13-scorecard}

\begin{scorecard}
\MT{lem:rg-formation-carrier} & Gauge-invariant residuals; triangle energy splits into contrasts and flux; orthogonal split into normal, loop, gauge directions & \SOLID & Standard in substance \\
\MT{prop:rg-formation-laws} & Normalized joint, quadratic restriction cost, free energy as relative entropy to the optimizer, gap \(O(N\mathrm T^2/\mathrm J_1)\); with \mtc{prop:rg-formation-normalized-experiment} & \SOLID & Standard in substance \\
\MT{prop:rg-formation-parent-potential} & Parents are exact marginals with Schur stiffness and bounded tilt; free energy drops by exactly \(\mathrm T\Delta\); with \mtc{prop:rg-formation-rate-composition} and \mtc{prop:rg-formation-score-dissipation}, reversible inherited rows and an H-theorem & \SOLID & Standard in substance: Gaussian integration, chain rule \\
\MT{def:rg-formation-kinetics} & Langevin plus heat-bath label jumps; declared parent kinetics; own-law learning (\mtc{def:rg-formation-learning}) & \CHOICE & The parent kinetics is a Markov closure \\
\MT{def:rg-formation-regime} & Doubling recursion, amplitude, temperature \(\epsilon^{20N-34}\), horizon, prepared mean, calibration & \CHOICE & Where the outcome is put in \\
\MT{thm:finite-depth-adaptive-formation}(b) & Component counts \(3\cdot2^{L-1}\to\cdots\to1\) on the windows & \TOY & Triangle chain, \((\R,+)\), scalar fibers, supplied hierarchy; linear relaxation; constants depend on \(L\) \\
\MT{thm:finite-depth-adaptive-formation}(a), (d) & Learning follows calibration; attention mass and hidden kernels move after formation & \TOY & Movement engineered; no effect on formation \\
\MT{thm:finite-depth-adaptive-formation}(c), (e) & Discarded loss and score error vanish; histories couple; paths concentrate & \TOY & Resolution \(\sqrt{\mathrm T}/\epsilon\), not total variation of path laws \\
\MT{rem:rg-formation-saturation} & Record saturated, learning calibration-dominated, loops thermal; coherence unchanged without record and learning & \SOLID & Negative content fixing the theorem's reading \\
\MT{prop:rg-formation-gaussian-core} & Exact loss; two-sided asymptotic; base-\(4/3\) separation criterion & \TOY & Sharp, for the Gaussian core of the chain only \\
\MT{prop:rg-formation-endogenous-cycle} & A pair forms by attention feedback with probability at least \(4.6\times10^{-3}\); exact parent over one time unit & \TOY & Four agents; gating, seed, circulation, \(\mathrm T=\epsilon^8\), no learning; symmetric law \\
\MT{prop:rg-gated-depth-one} & Stationary blocks are minimum-edge-cover components; each merge costs \(\mathrm T\) & \SOLID & No-go inside the gated family only; stationary; \(\mathrm T\to0\) at fixed graph \\
\MT{open:rg-formation-regime-sharpness} & Is the doubling needed; depth-uniform constants & \OPEN & Chain family, base-\(4/3\) schedules \\
\MT{open:rg-formation-adaptive-family} & Learning or holonomy changes, in law, which block forms & \OPEN & Four-cycle with learned rows \\
\MT{open:rg-formation-general} & A recursive (R14) second coalescence that is an episode (R15) & \OPEN & Outside the gated family \\
\end{scorecard}

\begin{hygiene}
Conditions that decide the chain's outcome, several chosen precisely to switch off the feedback the
program is interested in: the supplied strengths with the doubling, the prepared displacement with a strict
margin, the tie \(\mathrm T=\mathrm A^2\epsilon^6\), the calibration record, the saturated record, and source
palettes confined to each triangle. Conditions that exclude the program's targets: the group \((\R,+)\) with
scalar fibers, unit mobility (models as fast as beliefs), and the uncomputed \(\epsilon_L\) with
depth-dependent constants. For the four-cycle the conditions that matter are the supplied rows of the table
in \cref{comp:ch13-cycle}; for the gated family they are exclusive single-label attention, a scalar
translation group, and temperature-independent energies and prior ratios.
Hygiene: the root observation and gauge weight, the cancellation term \(-\mathrm TR_{\mathrm{ev}}\), the
preparation neighborhood of smooth transports, positivity and non-explosion of densities, the choice of
entrance constant and tolerance above their thresholds, the Gaussian reference, and \(\mathrm T\le1\).
\end{hygiene}

\section{Toy versus general}
\label{comp:ch13-toy}

None of the three results has hypotheses that include the program's target, 512 or more agents with
nonabelian, noncompact groups. For the triangle chain the mechanism-level content concerns coarse-graining,
not formation: exact marked parents can be carried through a hierarchical relaxation with separately
vanishing errors, provided the stiffness separation beats the adiabatic lag in thermal units; that this
tension appears in every construction of the type is plausible but unproved. The doubling, the exponents
20, 6, and 14, the constants \(10^{-4}\), \([1,21/10]\), \(1/200\), and \(57/100\), the sigmoid prior, and the
chain topology are specifics, and the theorem, though valid for arbitrarily many agents, is not a
large-\(N\) result. For the four-cycle the mechanism-level content is the attention-contraction feedback,
with an order-one flux deciding which edges stretch; it is local and could operate in large graphs, but
whether it forms many pairs at once, how neighboring pairs interact, and whether the law picks them are not
addressed, and the floor \(3e^{-1/4}/512\), the rates 8 and 4, and the one-unit horizon are specifics. For
the gated family the mechanism-level content, that exclusive gating puts the label into a sheet's weight and
never its stiffness, is proved on every finite graph; it is the chapter's one result that is not about a
small network, and it does not reach non-exclusive or multi-head attention, temperature-dependent
couplings, interactions that stiffen the directions between blocks, nonabelian links, dynamics, or many
agents at fixed temperature.

On the reader's own intuitions the chapter is neutral. Formation in the chain is continuous: the dynamics is
smooth relaxation, and the discrete counts come from the resolution \(\eta_L\mathrm A\) and scale separation,
not from a transition. No phase transition, condensation, or thermodynamic limit appears; every result is at
fixed finite \(N\) with \(\mathrm T\) or \(\epsilon\) tending to zero. The gated no-go shows fragmentation at fixed
\(N\) and low temperature in one family, and rules condensation neither in nor out elsewhere or as
\(N\to\infty\) at fixed temperature.

\section{Open ends}
\label{comp:ch13-open}

The scope paragraph's declined claims map onto three open problems. Order-one or nonabelian holonomy, a
model-channel record, slow per-agent models, and a record informative at the coherence scale are absent
from the chain and reappear as the ingredient list of \MT{open:rg-formation-adaptive-family}: a family in
which learning or holonomy changes, in law, which components form needs at least one of a record
standardized at the coherence scale, for example through \(\chi(r/(\eta_L\mathrm A))\); model means and links
both slow relative to beliefs, since a fast link or fast model erases the model disagreement on the belief
clock (\MT{prop:obs-slow-model-link-dichotomy}; see \cref{comp:ch08}); loops of order \(\mathrm A\) or
nonabelian links; or interacting model-channel records. The problem is posed on the four-cycle with learned
attention rows, as a total-variation difference, bounded below, between the conditional laws of the first
formed block in the learned and frozen arms, or between initial means differing only in their loop
coordinate, or a proof that these differences vanish for almost every mean. In a slower-mobility regime the
learned and frozen first-block laws coincide in the limit, so learning selects nothing there
(\MT{prop:phil-early-block-null}), and a parent activated after hidden thermalization has separately
vanishing errors on the seeded sub-event (\MT{prop:phil-early-seeded-parent}); the original coupled mobility
and the loop comparison remain open (\cref{comp:ch14}).

Depth-uniform constants and a non-doubling regime are \MT{open:rg-formation-regime-sharpness}; a positive
answer would still leave a stiffness range growing like a power of \(N\). Arbitrary families and initial laws,
mixed-level source support, permanent membership, and control as the depth tends to infinity are folded into
\MT{open:rg-formation-general}: a second endogenous coalescence, with learned parameters and from
non-matching assignments, in a source-gated model with temperature-independent data outside the gated
family, recursive in the sense of R14, held in the sense of R15, and followed by parents with separately
vanishing errors, or a proof that none exists in the class. A dynamic second coalescence inside the gated
family falls between this problem and the depth-one proposition. Intrinsic selection of the contrasts is
\MT{open:cg-natural-selector}, a persistent multilevel hierarchy formed by descent is
\MT{open:full-graph-hierarchy}, and behind all of them stands \MT{open:rg-pointwise-parent-dynamics}: the
chapter's parents are declared Markov closures checked against the fine process in two special families,
which does not settle how a parent's dynamics follows from the network it summarizes.

\begin{verdictbox}{\NARROWER}
\textbf{What the chapter says it does.} Its opening asks whether coherent components can form in succession
in one normalized experiment while every parent is activated from the complete preceding law, keeps every
loop and label, and learns its own parameter, and says two witnesses answer yes under explicit restrictions.
The opening paragraph of the introduction describes a finite-depth family that forms successive meta-agents
under free-energy descent with learning, and the program aim is that collectives form through the agents'
coupled inference. The introduction's reading map, by contrast, describes the chapter accurately.

\textbf{What the statements give.} For the chain the question as posed is answered yes
(\MT{thm:finite-depth-adaptive-formation}): exact marginal parents that keep loops and labels and learn from
their own laws track the fine process with separately vanishing errors at every finite depth. The succession
itself is linear relaxation of a supplied displacement through a supplied stiffness hierarchy; attention is
within \(10^{-4}\) of its prior, learning follows calibration and moves after formation, loops are thermal,
and the coherence proof works with attention and learning switched off (\MT{rem:rg-formation-saturation}).
The meta-agents are marginal laws, not coarse agents; the temperature is \(\epsilon^{20N-34}\) and the
constants depend on the depth without control. The four-cycle witness (\MT{prop:rg-formation-endogenous-cycle})
and the depth-one no-go (\MT{prop:rg-gated-depth-one}) are delivered as stated.

\textbf{The gap.} Nothing shows a collective forming because of coupled inference beyond one level. The one
endogenous witness is a sampled-path event on four agents with a symmetric law and no learning, whose end
state equilibrium already favors, and the opening's ``two witnesses answer yes'' fits the chain but not the
four-cycle, which has neither succession nor learning. Recursive endogenous formation
(\MT{open:rg-formation-general}), law-level selection by learning or holonomy
(\MT{open:rg-formation-adaptive-family}), and the depth-uniform control that any statement about hundreds of
agents needs (\MT{open:rg-formation-regime-sharpness}) are open. The technical chapter's scope paragraph is
candid about all of this; it is the introduction's opening summary and the program aim that promise more
than these theorems deliver.
\end{verdictbox}
